\documentclass[11pt]{article}
\usepackage[T1]{fontenc}
\usepackage{lmodern}
\usepackage[margin=1.05in]{geometry}
\usepackage{amsmath,amssymb,amsthm}
\usepackage{microtype}
\usepackage[hidelinks]{hyperref}
\newtheorem{theorem}{Theorem}[section]
\newtheorem{lemma}[theorem]{Lemma}
\newtheorem{proposition}[theorem]{Proposition}
\newtheorem{corollary}[theorem]{Corollary}
\theoremstyle{definition}
\newtheorem{definition}[theorem]{Definition}
\theoremstyle{remark}
\newtheorem{remark}[theorem]{Remark}
\newcommand{\R}{\mathbb R}
\newcommand{\one}{\mathbf 1}
\newcommand{\norm}[1]{\left\|#1\right\|}
\newcommand{\abs}[1]{\left|#1\right|}
\newcommand{\eps}{\varepsilon}
\newcommand{\intr}{\operatorname{int}}
\DeclareMathOperator{\spanR}{span}
\DeclareMathOperator{\sgn}{sgn}
\numberwithin{equation}{section}
\title{Robust permanence for endotactic power-law systems}
\author{Internal research reconstruction}
\date{7 October 2026}
\begin{document}
\maketitle
\begin{abstract}
We give a self-contained proof candidate for permanence of every finite
endotactic power-law edge system under arbitrary measurable reaction
rates in a supplied compact positive interval. Each positive
stoichiometric class has a compact convex set that is forward invariant
and absorbing for all admissible rate paths, with a path-independent
finite entry bound for each initial state. The proof reconstructs,
with attribution, the affine approximation lemma of the pinned OpenAI
release149 manuscript, and replaces that manuscript's weakly reversible
cut argument by a weighted essential-source comparison. The same
construction permits arbitrary measurable switching and convex mixing
among finitely many endotactic diagrams with a common stoichiometric
span. A separate invariance corollary allows arbitrary finite
compartment topologies with common transport coefficients across
species. This is an internally audited reconstruction; external
correctness, formal verification, and priority are not asserted.
\end{abstract}

\section{Interface, status, and main statement}
\label{sec:interface}

Let $d\geq1$. A finite power-law edge diagram has a finite vertex set
$C\subset\R^d$ and a finite set of directed, possibly labelled edges
$R$. Write an edge $e$ as $y_e\to y'_e$, with $y_e,y'_e\in C$, and put
$\nu_e=y'_e-y_e$. Self-loops and repeated labelled edges are harmless.
For $x\in\R_{>0}^d$ define
\[
 x^y=\prod_{i=1}^d x_i^{y_i},\qquad
 S=\spanR_{\R}\{\nu_e:e\in R\},\qquad
 P=(c+S)\cap\R_{>0}^d,
\]
where $c\in\R_{>0}^d$ is a fixed representative. A chemical reaction
network is the subcase $C\subset\mathbb Z_{\geq0}^d$. Real, including
negative, exponents are permitted in this manuscript, always on the
positive domain. No boundedness of $P$ is assumed.

\begin{definition}[Endotacticity]
\label{def:endotactic}
For $w\in\R^d$ let
\[
 R_w=\{e\in R:w\cdot\nu_e\ne0\}.
\]
If $R_w$ is nonempty, its essential support is the set of sources
maximizing $w\cdot y_e$ among $e\in R_w$. The diagram is
\emph{endotactic} if, for every $w$, every essential edge $e\in R_w$
whose source belongs to this support satisfies $w\cdot\nu_e<0$.
When $R_w$ is empty the
condition is vacuous.
\end{definition}

This is the essential-source convention of
\cite[Definitions 3.9 and 3.14]{GMS2014}. In particular the support is
not taken among all sources or all vertices. Strong endotacticity
imposes an additional condition at extrema among all sources and is
not assumed here. For a fixed vector $r$ not in $S^\perp$, applying
Definition~\ref{def:endotactic} to $-r$ says that every essential edge
at the \emph{minimum} $r$-source level has $r\cdot\nu_e>0$.

Fix supplied constants $0<\kappa\leq K<\infty$. Rates
$k_e:[0,\infty)\to[\kappa,K]$ may be arbitrary measurable functions,
and the equation is
\begin{equation}
 \dot x(t)=f(t,x(t)):=\sum_{e\in R}k_e(t)x(t)^{y_e}\nu_e.
 \label{eq:system}
\end{equation}
Changing a rate on a null set does not change a solution. A solution
means an absolutely continuous Carath\'eodory solution satisfying
\eqref{eq:system} almost everywhere. The vector field is measurable
in time and uniformly locally Lipschitz in state on the positive
domain. Standard Carath\'eodory local existence and uniqueness
therefore apply. A positive class is preserved because
$f(t,x)\in S$ almost everywhere.

\begin{theorem}[A common classwise absorber]
\label{thm:main}
For every finite endotactic diagram, supplied interval $[\kappa,K]$,
and positive class $P$, there is a nonempty compact convex set
$\mathcal K_P\subset P$ with the following properties. Every
measurable rate path in the interval and every $x^0\in P$ give a
unique global positive solution of \eqref{eq:system}. It enters
$\mathcal K_P$ in finite time and remains there. The set is common
to all admissible rate paths. For each $x^0$, its entry time has a
finite bound independent of the rate path. In particular there is
$\epsilon_P\in(0,1)$, independent of the rate path and initial
point, such that every solution eventually satisfies
\[
 \epsilon_P\leq x_i(t)\leq\epsilon_P^{-1}\qquad(1\leq i\leq d).
\]
\end{theorem}

The proof of the theorem occupies Sections~\ref{sec:affine}--
\ref{sec:absorption}. Section~\ref{sec:switching} states the precise
common-span extension. Standard local ODE existence, elementary
absolute-continuity facts, compactness, and (only for the equilibrium
and periodic corollaries) Brouwer's theorem are the analytical
dependencies. The nonstandard affine lemma is proved in full below.

\paragraph{Provenance and research status.}
The affine construction is reconstructed from
\cite[Lemma 2.1 and Corollary 2.2]{OAI149}, at the pinned release
commit stated in the bibliography. That source itself claims
fixed-positive-rate permanence for weakly reversible chemical
networks; it does not state the all-endotactic measurable-rate
theorem above. The new deduction in this manuscript is the weighted
essential-source comparison, its path-independent strict drift,
and its transfer through the plateaux to common absorption.
All proof details have been internally reconstructed, but this
document makes no assertion of independent external acceptance or
completed formal verification. A targeted primary-source review is
summarized in Section~\ref{sec:scope}; absence of a found competing
theorem is not a novelty certificate.

\section{Finite affine minima with arbitrary positive tolerance}
\label{sec:affine}

For $m\geq1$, $u\in\R^m$, and $0<h<1$, let
\[
 X_h(u)=(h^{u_1},\ldots,h^{u_m}),\qquad
 \mu(u)=\min_i u_i.
\]
A label is a slope $r\in\R^m$, fixed independently of $h$, together
with an offset $b(h)$ defined for all sufficiently small positive
$h$. For a finite nonempty label family $\Lambda$ put
\[
 F_h(x)=\min_{\lambda\in\Lambda}
          (r_\lambda\cdot x+b_\lambda(h)).
\]
A label is \emph{active} at $x$ if it attains this minimum. Every
minimizer is active at a tie. In any finite family the offsets are
restricted to a common small interval of definition.

\begin{lemma}[Affine approximation; reconstructed from \cite{OAI149}]
\label{lem:affine}
Let $a<b$ be real, $m\geq1$, and let
$E:[a,b]^m\to(0,\infty)$ be any function. There is a finite
nonempty label family with $r_\lambda\in[a,b]^m$ and
$b_\lambda(h)=o(h^a)$ such that the following holds. If
$h_n\to0+$, $p^n\in\R^m$ tends to $p\in[a,b]^m$, and one fixed
label $\lambda$ is active at $X_{h_n}(p^n)$ for every $n$, then
\[
 \norm{p-r_\lambda}_\infty<E(r_\lambda).
\]
Neither continuity of $E$ nor membership $p^n\in[a,b]^m$ is
required. The family and slopes are fixed before $h$ is chosen.
\end{lemma}

\begin{proof}
We use induction on $m$. Throughout, exponential comparisons are
made only with strictly different \emph{limiting} exponents: if
$p_i>q$ and $p^n\to p$, then $h_n^{p_i^n}=o(h_n^q)$, whereas
if $p_i<q$, the ratio $h_n^{p_i^n}/h_n^q$ tends to infinity.
These facts follow by bounding $p_i^n-q$ away from zero with the
appropriate sign.

\medskip
\noindent\emph{A local family on the slice $\mu(p)=t$.}
For each $t\in[a,b]$ we construct a finite family $\Lambda_t$
containing $(t\one,0)$, all of whose slopes have minimum entry $t$,
and which has the lemma's sequential property when $\mu(p)=t$.
For $m=1$, or for $t=b$, the singleton $(t\one,0)$ suffices.
This supplies the local base case; the gluing argument below will
also prove the full statement in dimension one.

Suppose $m>1$ and $t<b$. Call a label type $k$ when exactly $k$
coordinates of its slope exceed $t$. Construct types
$k=1,\ldots,m-1$ successively, after the type-zero label
$(t\one,0)$. Once type $k-1$ has been constructed, choose
\begin{equation}
 t<\beta_k<b,\qquad
 \beta_k<\beta_{k-1}\ (k>1),\qquad
 \beta_k-t<E(r)\quad\text{for every type-$(k-1)$ slope }r.
 \label{eq:beta}
\end{equation}
This is possible because there are finitely many preceding slopes
and each tolerance there is positive.

For each $k$-element coordinate set $U$ apply the induction
hypothesis in dimension $k<m$ to $[\beta_k,b]^U$, with tolerance
$s\mapsto E(r(s))$, where $r(s)$ equals $s$ on $U$ and $t$ on
its complement. Every resulting lower-dimensional label
$(s,d(h))$ is lifted to
\begin{equation}
 r=r(s),\qquad c(h)=-h^{\beta_k}+d(h),\qquad d(h)=o(h^{\beta_k}).
 \label{eq:lift}
\end{equation}
The family remains finite. Every lifted label has exactly $k$
raised entries and offset $o(h^t)$.

To verify the local property, fix a label $(r,c)$ of type $k$
that stays active along $h_n\to0$, $p^n\to p\in[a,b]^m$ with
$\mu(p)=t$. Let $U$ be its raised coordinate set and write $h=h_n$.
For $k>0$, comparison with $(t\one,0)$ gives
\[
 \sum_{i\in U}(r_i-t)h^{p_i^n}
       \leq h^{\beta_k}-d(h)=(1+o(1))h^{\beta_k}.
\]
All coefficients on the left are at least $\beta_k-t>0$.
A limiting coordinate $p_i<\beta_k$ would make its term divided
by $h^{\beta_k}$ diverge, so
\begin{equation}
 p_i\geq\beta_k\quad(i\in U).
 \label{eq:raised}
\end{equation}
Among labels lifted on this same $U$, the terms off $U$ and
$-h^{\beta_k}$ are identical. Thus $(s,d)$ is active within the
lower-dimensional family. Its limiting exponent vector belongs
to $[\beta_k,b]^U$, so the induction hypothesis yields
\[
 \max_{i\in U}|p_i-r_i|<E(r).
\]
The vectors $p_U^n$ may approach a face of this cube from outside;
the induction statement deliberately permits this.

Now suppose $k<m-1$. We claim $p_i\leq\beta_{k+1}$ for every
$i\notin U$, including the type-zero case $U=\varnothing$.
Otherwise choose an unraised $j$ with $p_j>\beta_{k+1}$ and any
type-$(k+1)$ label lifted on $U\cup\{j\}$. On this set every
limiting exponent is strictly above $\beta_{k+1}$, by
\eqref{eq:raised}, $\beta_k>\beta_{k+1}$, and the choice of $j$.
The competitor's value minus that of $(t\one,0)$ is therefore
\[
 -h^{\beta_{k+1}}+o(h^{\beta_{k+1}}).
\]
The supposed active label's excess over $(t\one,0)$ is zero
when $k=0$, and is at least $-h^{\beta_k}+o(h^{\beta_k})$
otherwise. Since $h^{\beta_k}=o(h^{\beta_{k+1}})$, the competitor
is strictly smaller for sufficiently small $h$ along the sequence,
a contradiction. Consequently every unraised coordinate satisfies
\[
 |p_i-r_i|=p_i-t\leq\beta_{k+1}-t<E(r).
\]
If $k=m-1$, the single unraised coordinate instead equals $t$,
because $\mu(p)=t$ and all raised coordinates have limits strictly
above $t$. Its error is zero. This proves the local property for
all types.

\medskip
\noindent\emph{A neighborhood of the slice.}
For a fixed local label $\lambda$, let $A_\lambda\subset[a,b]^m$
be the set of limits of sequences $p^n\in\R^m$ along which
$h_n\to0$ and that label is active in $\Lambda_t$. This set is
closed. Indeed, if $z_j\in A_\lambda$ tends to $z$, select from a
sequence realizing $z_j$ one point $u_j$ and scale $0<h_j<1/j$
with $\norm{u_j-z_j}_\infty<1/j$ and actual activity. The selected
points tend to $z$ and realize $z\in A_\lambda$.

The bad limiting set
\[
 B_t=\bigcup_{\lambda=(r,c)\in\Lambda_t}
      \{p\in A_\lambda:\norm{p-r}_\infty\geq E(r)\}
\]
is compact. Each $E(r)$ here is a fixed real number and the union
is finite; no continuity of $E$ is needed. The local property
shows that $B_t$ contains no point with $\mu(p)=t$. By continuity
of $\mu$ and compactness there is $\rho_t>0$ such that no point
of $B_t$ has minimum exponent in
\[
 I_t=(t-\rho_t,t+\rho_t).
\]
If $B_t$ is empty any positive radius first suffices. Otherwise
choose $\rho_t<\min_{p\in B_t}|\mu(p)-t|$.
Shrink this radius further so that every local offset satisfies
\begin{equation}
 c(h)=o(h^{t+\rho_t}).
 \label{eq:offsetdecay}
\end{equation}
For a nonsingleton family it suffices to choose
$t+\rho_t<\min_k\beta_k$ in \eqref{eq:lift}; singleton offsets
are zero. The local family's sequential guarantee now holds
whenever the limiting minimum belongs to $I_t$.

\medskip
\noindent\emph{Finite gluing.}
Choose a finite subcover of $[a,b]$ by the $I_t$. Remove every
interval contained in another, retaining one in case of equality,
and order the remaining centers $t_1<\cdots<t_N$. Write
$I_{t_j}=(L_j,R_j)$. Noncontainment forces both lists of endpoints
to be strictly increasing in this center order. The first interval
contains $a$ and the last contains $b$: for example $L_1<a$ because
some interval covers $a$, and $R_1>t_1\geq a$.
Consecutive intervals overlap. Otherwise $R_j\leq L_{j+1}$,
and the point $R_j\in(t_j,t_{j+1})\subset[a,b]$ lies in none
of these open intervals.

For $1\leq j<N$ choose a cut
\begin{equation}
 \max\{t_j,L_{j+1}\}<q_j<\min\{t_{j+1},R_j\}.
 \label{eq:cutchoice}
\end{equation}
This interval is nonempty because consecutive intervals overlap
and each contains its center. With $q_0=a$, $q_N=b$, it follows
that
\begin{equation}
 t_j<q_j<t_{j+1}\quad(j<N),\qquad
 [q_{j-1},q_j]\subset I_{t_j}\quad(1\leq j\leq N).
 \label{eq:cutranges}
\end{equation}
If $N=1$ the one interval covers $[a,b]$ and no cuts are required.

Combine all selected local families, replacing each offset in
family $j$ by
\begin{equation}
 \widetilde c(h)=c(h)-\sum_{\ell<j}h^{q_\ell}.
 \label{eq:shifts}
\end{equation}
Every resulting offset is $o(h^a)$: in \eqref{eq:offsetdecay}
the exponent $t_j+\rho_{t_j}$ exceeds $a$, and every shift
exponent $q_\ell>t_\ell\geq a$. A common shift preserves activity
within its local family.

Suppose a fixed label in family $j$ is globally active along
$h_n\to0$, $p^n\to p\in[a,b]^m$. We show
$q_{j-1}\leq\mu(p)\leq q_j$. If $j<N$ and $\mu(p)>q_j$,
subtract the shifted value of $(t_{j+1}\one,0)$ from its value.
The difference is
\[
 (r-t_{j+1}\one)\cdot X_h(p^n)+c(h)+h^{q_j}
             =h^{q_j}+o(h^{q_j})>0,
\]
because all limiting coordinate exponents exceed $q_j$, and
$c(h)=o(h^{q_j})$ by $q_j<R_j=t_j+\rho_{t_j}$. This contradicts
activity. If $j>1$ and $\mu(p)<q_{j-1}$, comparison with shifted
$(t_{j-1}\one,0)$ instead gives
\[
 (r-t_{j-1}\one)\cdot X_h(p^n)+c(h)-h^{q_{j-1}}.
\]
Every coefficient in the scalar product is at least
$t_j-t_{j-1}>0$, and some coordinate exponent has limit below
$q_{j-1}$. Dividing by $h^{q_{j-1}}$ therefore gives $+\infty$,
since $c(h)=o(h^{q_{j-1}})$ by $q_{j-1}<t_j$. Again activity
is contradicted. The outer endpoint bounds follow from
$\mu(p)\in[a,b]$.

By \eqref{eq:cutranges}, including at either cut, $\mu(p)$ belongs
to $I_{t_j}$. Global activity implies activity within the
unshifted local family, so its neighborhood guarantee gives
$\norm{p-r}_\infty<E(r)$. This finishes the induction.
\end{proof}

\begin{corollary}[Uniform activity at a fixed finite family]
\label{cor:activity}
For a family supplied by Lemma~\ref{lem:affine}, there is
$h_{\rm act}\in(0,1)$ such that, for every $0<h<h_{\rm act}$,
$p\in[a,b]^m$, and label active at $X_h(p)$,
\[
 \norm{p-r_\lambda}_\infty<E(r_\lambda).
\]
\end{corollary}
\begin{proof}
If no such cutoff existed, take $h_n\to0$ and violations at
$p^n\in[a,b]^m$. Finiteness yields a subsequence with one fixed
active label; compactness yields a further subsequence with
$p^n\to p$. The violating weak inequality passes to the limit
because $E$ is evaluated at that fixed slope. This contradicts
Lemma~\ref{lem:affine}.
\end{proof}

\begin{remark}
The proof handles ties, negative $a$, the slice $t=b$, and
approaches to lower-dimensional faces from outside. Its offsets
can in fact be taken as finite negative sums of powers
$-\sum_\ell h^{\gamma_\ell}$ with $\gamma_\ell>a$; this follows
from the lifting and shifting operations. No useful bound on the
number of labels or the numerical cutoff is established here.
\end{remark}

\section{Active-slope comparisons and weighted endotactic drift}
\label{sec:flux}

Apply the affine construction with $m=d$, $a=-1$, $b=1$. Let
\begin{equation}
 D=\{e_i:1\leq i\leq d\}\cup
       \{z-y:y,z\in C,\ z\ne y\},
 \label{eq:D}
\end{equation}
a finite set of nonzero vectors. For $r\in[-1,1]^d$ define
\begin{equation}
 E(r)=\min\left(\{1/2\}\cup
 \left\{\frac{|r\cdot v|}{2\norm{v}_1}:
             v\in D,\ r\cdot v\ne0\right\}\right)>0.
 \label{eq:E}
\end{equation}
Zero products are omitted. The function can be discontinuous
on comparison hyperplanes, which is why the preceding lemma's
pointwise-positive interface matters. Fix its finite family and
write
\[
 L_{\lambda,h}(x)=r_\lambda\cdot x+b_\lambda(h),\qquad
 F_h(x)=\min_\lambda L_{\lambda,h}(x),\qquad
 A_h(x)=\{\lambda:L_{\lambda,h}(x)=F_h(x)\}.
\]
Let $B_h=[h,h^{-1}]^d$. Every $x\in B_h$ is $X_h(p)$ for the
unique $p=\log x/\log h\in[-1,1]^d$. For small $h$, activity
gives $\norm{p-r}_\infty<E(r)$ for every active slope $r$.
Only after fixing this family define
\begin{equation}
 H=\min\left(\{1\}\cup
 \{|r_\lambda\cdot v|/2:\lambda\in\Lambda,\ v\in D,
                                      \ r_\lambda\cdot v\ne0\}\right)>0.
 \label{eq:H}
\end{equation}
Finiteness is essential to this positive minimum. For $v\in D$
with $r\cdot v\ne0$,
\[
 |(p-r)\cdot v|\leq\norm{p-r}_\infty\norm{v}_1
       < E(r)\norm{v}_1\leq |r\cdot v|/2.
\]
Consequently
\begin{equation}
 \sgn(p\cdot v)=\sgn(r\cdot v),\qquad |p\cdot v|\geq H.
 \label{eq:sign}
\end{equation}
For two sources $u,w$ with $r\cdot u>r\cdot w$ this implies
\begin{equation}
 \frac{x^u}{x^w}=h^{p\cdot(u-w)}\leq h^H.
 \label{eq:ratio}
\end{equation}
The comparison includes sources in different graph components.
It also includes every nonzero reaction vector, even when its
target is not a source. Neither all vertices being sources nor
weak reversibility has been assumed.

Put
\begin{equation}
 L=\sum_{e\in R}\norm{\nu_e}_1,\qquad
 M=\max_{e\in R}\norm{y_e}_1,
 \label{eq:LM}
\end{equation}
with $M=0$ for an empty edge set. The use of an absolute norm
allows real exponents as well as chemical complexes.

\begin{lemma}[Weighted essential-source domination]
\label{lem:drift}
Suppose $h$ is sufficiently small for activity and satisfies
\begin{equation}
 K L h^H\leq\kappa H.
 \label{eq:smallh}
\end{equation}
For every $x\in B_h$, every active slope $r$, and every choice of
rate values $k_e\in[\kappa,K]$,
\[
 r\cdot\sum_e k_e x^{y_e}\nu_e\geq0.
\]
If $r\notin S^\perp$, then the stronger uniform bound holds:
\begin{equation}
 r\cdot\sum_e k_e x^{y_e}\nu_e
       \geq\delta_h:=\kappa Hh^M>0.
 \label{eq:drift}
\end{equation}
\end{lemma}

\begin{proof}
For $r\in S^\perp$ every summand has zero scalar product.
Otherwise $R_{-r}$ is nonempty. Choose an edge
$e_*:w\to w'$ minimizing $r\cdot y_e$ among $e\in R_{-r}$.
Endotacticity in direction $-r$ says that every essential edge
whose source is at this minimum level has $r\cdot\nu_e>0$.
In particular, $r\cdot\nu_{e_*}>0$. Since the nonzero vector
$\nu_{e_*}$ belongs to $D$, \eqref{eq:H} gives
$r\cdot\nu_{e_*}\geq2H$.

Every edge with $r\cdot\nu_e<0$ is essential. Its source
$u=y_e$ is strictly above the minimum level, since equality
would contradict the preceding endotactic condition. Thus
$x^{y_e}\leq h^H x^w$ for all negative summands by
\eqref{eq:ratio}. Also $\norm{r}_\infty\leq1$, so
$|r\cdot\nu_e|\leq\norm{\nu_e}_1$. Discarding all positive
terms except the selected one yields
\begin{align}
 r\cdot\sum_e k_e x^{y_e}\nu_e
 &\geq\kappa x^w(r\cdot\nu_{e_*})
       -K\sum_{r\cdot\nu_e<0}x^{y_e}|r\cdot\nu_e| \notag\\
 &\geq x^w(2\kappa H-KLh^H)
 \geq\kappa Hx^w.\label{eq:weighted}
\end{align}
For $x=X_h(p)\in B_h$, $p\cdot w\leq\norm{w}_1\leq M$.
Because $0<h<1$,
$x^w=h^{p\cdot w}\geq h^M$. This proves \eqref{eq:drift}.
\end{proof}

\begin{remark}
The factors $r\cdot\nu_e$ cannot be ignored. A comparison of
unweighted reaction fluxes alone would not justify the derivative
sign for a general endotactic diagram. The estimate above controls
the entire negative sum using one positive edge and a finite
coefficient budget. It holds for every point of the rate box,
so the scale and drift bound are independent of rate switching.
\end{remark}

\section{Compact plateaux and absolutely continuous trajectories}
\label{sec:plateau}

For a fixed positive representative $c$ choose $h_P\in(0,1)$
so small that every $0<h\leq h_P$ is in the common definition
interval of the offsets, satisfies the activity bound and
\eqref{eq:smallh}, and obeys
\begin{equation}
 \norm{\frac{\log c}{\log h}}_\infty<H.
 \label{eq:cscale}
\end{equation}
Such a choice is possible: the family and $H$ are fixed, $h^H$
tends to zero, and $1/|\log h|$ tends to zero. In particular
$c\in\intr B_h$ for these scales. Define
\[
 M_h=F_h(c),\qquad \mathcal K_h=\{x\in\R^d:F_h(x)=M_h\}.
\]

\begin{proposition}[The global compact plateau]
\label{prop:plateau}
For $0<h\leq h_P$, one label has zero slope and constant value
$M_h$, and $F_h\leq M_h$ on all of $\R^d$. The plateau
$\mathcal K_h$ is nonempty, compact, and convex, and
\begin{equation}
 \mathcal K_h\subset\intr B_{\sqrt h}.
 \label{eq:innerbox}
\end{equation}
\end{proposition}
\begin{proof}
At $c$, any nonzero coordinate of an active slope would, by
\eqref{eq:sign} with $v=e_i$, force
$|\log c_i/\log h|\geq H$, contradicting \eqref{eq:cscale}.
Thus every active slope there is zero. There is at least one
active label, giving the constant value $M_h$ and the global
upper bound $F_h\leq M_h$. Its value is the least value of all
constant labels. Therefore
\begin{equation}
 \mathcal K_h=\bigcap_{\lambda\in\Lambda}
                 \{x:L_{\lambda,h}(x)\geq M_h\}
 \label{eq:polytope}
\end{equation}
is closed and convex, and contains $c$.

At every $x=X_h(p)\in\mathcal K_h\cap B_h$ the least constant
label is active. Its activity bound gives
$\norm{p}_\infty<E(0)=1/2$, so
$x\in\intr B_{\sqrt h}$. Thus the plateau misses the boundary
of $B_h$. If it contained a point outside $B_h$, the segment
from $c\in\intr B_h$ to that point would meet this boundary
inside the plateau by convexity, a contradiction. The whole
plateau therefore lies in $B_h$ and then in its inner box.
Closedness and boundedness prove compactness.
\end{proof}

\begin{lemma}[No neutral active slopes below the plateau in a class]
\label{lem:neutral}
If $x\in P\setminus\mathcal K_h$, every active slope lies outside
$S^\perp$.
\end{lemma}
\begin{proof}
For $r_\lambda\in S^\perp$, $x-c\in S$ implies
$L_{\lambda,h}(x)=L_{\lambda,h}(c)\geq M_h$.
But the global upper bound and $x\notin\mathcal K_h$ give
$F_h(x)<M_h$, so this label cannot be active.
\end{proof}

\begin{lemma}[Chain rule for a finite affine minimum]
\label{lem:chain}
If $x:[s,t]\to\R^d$ is absolutely continuous, so is
$F_h\circ x$, and for almost every $u\in[s,t]$,
\begin{equation}
 \frac{d}{du}F_h(x(u))
     =\min_{\lambda\in A_h(x(u))}r_\lambda\cdot\dot x(u).
 \label{eq:chain}
\end{equation}
\end{lemma}
\begin{proof}
Every affine composition is absolutely continuous. A finite
minimum of such functions is absolutely continuous, for example
by the identity $\min(a,b)=(a+b-|a-b|)/2$ and the preservation
of absolute continuity by the absolute-value map.
At a time $u$ where $x$ is differentiable, inactive labels have
positive gaps above the minimum. Finiteness and continuity
exclude them from the minimum for sufficiently small positive
time increments. Each active label has the expansion
\[
 L_{\lambda,h}(x(u+v))
   =F_h(x(u))+v r_\lambda\cdot\dot x(u)+o(v)
       \qquad(v\downarrow0).
\]
Taking the minimum gives the right derivative in
\eqref{eq:chain}. At almost every time, both $x$ and the
absolutely continuous composition are differentiable, so this
right derivative is the composition's derivative. All
minimizers at a tie are included; a tie with differing two-sided
derivatives is confined to the null set where the minimum is
not differentiable.
\end{proof}

\begin{proposition}[Uniform increase, invariance, and finite entry]
\label{prop:entry}
At any fixed scale $0<h\leq h_P$, the following hold for every
measurable admissible rate path.
\begin{enumerate}
\item On every solution segment contained in $B_h$, $F_h(x(t))$
is nondecreasing.
\item If a segment $x([s,t])$ is contained in
$B_h\cap(P\setminus\mathcal K_h)$, then
\begin{equation}
 F_h(x(t))-F_h(x(s))\geq\delta_h(t-s),
 \qquad \delta_h=\kappa Hh^M.
 \label{eq:increase}
\end{equation}
\item Every solution starting in $\mathcal K_h$ is global and
positive and stays in $\mathcal K_h$.
\item A global solution in $P$ staying in $B_h$ from time $t_0$
onward enters $\mathcal K_h$ no later than $t_0+W_h$, where
\begin{equation}
 W_h=\frac{M_h-\min_{x\in B_h}F_h(x)}{\kappa Hh^M}<\infty.
 \label{eq:Wh}
\end{equation}
\end{enumerate}
\end{proposition}
\begin{proof}
The first two assertions follow by applying
Lemma~\ref{lem:chain} almost everywhere, using
Lemma~\ref{lem:drift} and then Lemma~\ref{lem:neutral}, and
integrating. This is the appropriate replacement of an
all-times right-derivative argument for merely measurable rates.

For the third assertion, a maximal positive solution starting
in $\mathcal K_h\subset\intr B_h$ initially stays in $B_h$.
On every interval where it does, monotonicity and the global
upper bound $F_h\leq M_h$ force $F_h(x(t))=M_h$.
It cannot reach the boundary of $B_h$ while in its maximal
interval: continuity would place such a first exit point in
both $\mathcal K_h$ and this boundary, contrary to
Proposition~\ref{prop:plateau}. Thus the solution stays in the
compact positive set $\mathcal K_h$ throughout its maximal
interval. The vector field is uniformly bounded and uniformly
locally Lipschitz on a positive neighborhood of this set. At a
finite maximal endpoint the solution consequently has a limit
inside $\mathcal K_h$, and Carath\'eodory local existence from
that point continues it. Hence it is global and stays positive.
The almost-everywhere inclusion $f(t,x)\in S$ also preserves
its class.

For the last assertion, if the trajectory avoided the plateau
until $t_0+W_h$, integrating \eqref{eq:increase} would give
\[
 F_h(x(t_0+W_h))\geq F_h(x(t_0))+\delta_h W_h\geq M_h.
\]
The global upper bound then places the trajectory on the
plateau by that time. This also follows by a limiting argument
if the first entry occurs at the endpoint. The third assertion
keeps it there afterward. If $W_h=0$, the whole box already
has plateau value and the assertion is immediate.
\end{proof}

\section{A common absorber on an unbounded class}
\label{sec:absorption}

\begin{proof}[Proof of Theorem~\ref{thm:main}]
If $S=\{0\}$, every reaction vector is zero, $P=\{c\}$,
and $\mathcal K_P=\{c\}$ proves the theorem directly. Otherwise
use the preceding construction. The scale $h_P$ was chosen
from the diagram, interval, and representative of $P$ before
choosing an initial point or a rate path.

Let $x^0\in P$. Choose $0<h_0\leq h_P$ such that
\[
 \norm{\frac{\log x^0}{\log h_0}}_\infty<H.
\]
Then $x^0\in B_{h_0}$ and, by the coordinate comparisons, every
active label there has zero slope. An active constant must have
the least constant value, which is $M_{h_0}=F_{h_0}(c)$.
Thus $x^0\in\mathcal K_{h_0}$. Proposition~\ref{prop:entry}
already proves that its solution is global and positive and
stays in this initial-point-dependent compact set, for any
admissible rate path.

Define
\begin{equation}
 h_{j+1}=\min\{\sqrt{h_j},h_P\}.
 \label{eq:scalesequence}
\end{equation}
If a solution has entered $\mathcal K_{h_j}$ and stays there,
then \eqref{eq:innerbox} and $h_{j+1}\leq\sqrt{h_j}$ imply
that it stays in $B_{h_{j+1}}$. The last part of
Proposition~\ref{prop:entry} makes it enter
$\mathcal K_{h_{j+1}}$ within an additional $W_{h_{j+1}}$ units
of time, and that set is then invariant. Induction gives this
assertion at each scale. No nesting of the plateaux is used.

Repeated square roots of $h_0\in(0,1)$ tend to 1. Since
$h_P<1$, the capped sequence reaches $h_P$ after finitely many
transitions, say $N$. One can take
\[
 N=\begin{cases}
 0,&h_0=h_P,\\[2pt]
 \left\lceil\log_2\bigl(|\log h_0|/|\log h_P|\bigr)\right\rceil,
          &h_0<h_P.
 \end{cases}
\]
An entry bound for the common final plateau is
\begin{equation}
 T(x^0)\leq\sum_{j=1}^N W_{h_j}.
 \label{eq:entrybound}
\end{equation}
All scales and terms here are independent of the rate path.

Finally put $\mathcal K_P=\mathcal K_{h_P}\cap(c+S)$.
It contains $c$, is compact and convex, and is contained in
the positive class by \eqref{eq:innerbox}. Both the plateau
and the affine class are invariant, so it is forward invariant
and absorbing. Taking $\epsilon_P=h_P$ gives the asserted
coordinate bounds. This proves the theorem also when $P$ is
unbounded.
\end{proof}

\section{Switching, convex mixtures, and invariant transport}
\label{sec:switching}

Let $(C_j,R_j)$, $1\leq j\leq J$, be a finite nonempty family
of endotactic diagrams in the same space $\R^d$. For edges
of diagram $j$ write $y_{j,e}$ and $\nu_{j,e}$, and let
\[
 S_j=\spanR_{\R}\{\nu_{j,e}:e\in R_j\}.
\]
Each edge has a measurable rate $k_{j,e}(t)\in[\kappa,K]$.
Let $\alpha_j(t)\geq0$ be measurable with
$\sum_j\alpha_j(t)=1$, and consider
\begin{equation}
 \dot x(t)=\sum_{j=1}^J\alpha_j(t)f_j(t,x(t)),\qquad
 f_j(t,x)=\sum_{e\in R_j}k_{j,e}(t)x^{y_{j,e}}\nu_{j,e}.
 \label{eq:mixture}
\end{equation}
Pure switching is the case that exactly one weight is 1 at
almost every time. Neither a dwell time nor a positive lower
bound on each individual mixture weight is imposed.

\begin{theorem}[Common-span switching and convex mixing]
\label{thm:switching}
Suppose every diagram has the same stoichiometric subspace:
$S_j=S$ for every $j$. For each positive class
$P=(c+S)\cap\R_{>0}^d$, there is a nonempty compact convex
$\mathcal K_P\subset P$ that is invariant and absorbing for
every system \eqref{eq:mixture}, simultaneously for all
admissible measurable weights and rates. All solutions from
$P$ are unique, global, and positive. Each initial state has
a finite entry bound independent of both the rates and the
weight path.
\end{theorem}
\begin{proof}
If $S=\{0\}$ all fields are zero and the assertion is immediate.
Otherwise make one affine construction using $C=\bigcup_j C_j$
in \eqref{eq:D}. Choose $H$ from its fixed finite slopes and
use the uniform budgets
\begin{equation}
 L=\max_j\sum_{e\in R_j}\norm{\nu_{j,e}}_1>0,\qquad
 M=\max_{j,e\in R_j}\norm{y_{j,e}}_1.
 \label{eq:familybudget}
\end{equation}
For any active $r\notin S^\perp$ and \emph{each} diagram $j$,
the essential set for $-r$ is nonempty because $S_j=S$.
The proof of Lemma~\ref{lem:drift}, using that diagram's
minimum essential source and the common comparison family,
gives
\[
 r\cdot f_j(t,x)\geq\kappa Hh^M=\delta_h
             \qquad(x\in B_h).
\]
For $r\in S^\perp$ every component scalar product is zero.
Convex mixing preserves both statements:
\[
 r\cdot\sum_j\alpha_j(t)f_j(t,x)\geq\delta_h
                      \quad(r\notin S^\perp).
\]
All choices are made using the whole finite family and the
supplied rate interval. The common span preserves $P$ and
excludes neutral active slopes below the plateau exactly as
in Lemma~\ref{lem:neutral}. The plateau, almost-everywhere
chain rule, extension, finite entry, and scale iteration
proofs now apply unchanged. The mixture is still measurable
in time and uniformly locally Lipschitz in state. The same
waiting bounds \eqref{eq:Wh} and \eqref{eq:entrybound} use
the common budgets \eqref{eq:familybudget} and are independent
of the weight path.
\end{proof}

\begin{remark}[Different spans do not justify common absorption]
\label{rem:span}
Take the two chemical diagrams $0\rightleftarrows A$ and
$0\rightleftarrows B$, with all rates 1, in two species.
They are endotactic but have spans $\R e_1$ and $\R e_2$.
Their combined span is $\R^2$. Switch permanently to the
first diagram. Then $x_B(t)=x_B(0)$, which can be arbitrarily
large in the positive class for the combined span. No one
compact subset of that class can absorb every initial state
under every switching path. Thus replacing the common-span
hypothesis by a statement about the union span would be
incorrect. Common invariance for supplied bounded initial
ranges, addressed next, is a different assertion.
\end{remark}

\begin{corollary}[A common invariant set without a common span]
\label{cor:safety}
For any finite family of endotactic diagrams and any compact
$Q\subset\R_{>0}^d$, there is a compact convex
$\mathcal K\subset\R_{>0}^d$ containing $Q$ that is forward
invariant for all measurable convex mixtures
\eqref{eq:mixture}, even when the $S_j$ differ.
\end{corollary}
\begin{proof}
Use the union comparison family and budgets
\eqref{eq:familybudget}, with $M=0$ if there are no edges.
Choose a scale small enough for activity,
$KLh^H\leq\kappa H$ (automatic if $L=0$), and
\[
 \sup_{q\in Q\cup\{\one\}}
       \norm{\frac{\log q}{\log h}}_\infty<H.
\]
The supremum is finite before division by $|\log h|$, because
$Q$ is a positive compact set. Use the representative
$c=\one$ to define the least-constant plateau
$\mathcal K=\mathcal K_h$. The coordinate comparison shows
that every point of $Q$ has an active constant label and
therefore belongs to this plateau. It is compact, convex,
and positive by Proposition~\ref{prop:plateau}.

For any active slope and any diagram, the weighted flux
proof gives a nonnegative scalar product: it is zero if the
slope is orthogonal to that diagram's span, and strictly
positive otherwise. Thus every mixture has nonnegative
active-slope derivative. The chain rule and compact-plateau
argument prove invariance and global positivity. No
classwise strict drift or absorbing assertion is used.
\end{proof}

\begin{corollary}[Common-species transport on arbitrary finite topologies]
\label{cor:transport}
Fix a finite family of endotactic diagrams and a compact
positive initial range $Q$. Let $\mathcal K$ be the plateau
constructed in Corollary~\ref{cor:safety}. For any finite
number $n$ of compartments, let the state of compartment $z$
be $x_z\in\R_{>0}^d$ and consider
\begin{equation}
 \dot x_z=f_z(t,x_z)+\sum_{w=1}^n m_{zw}(t)(x_w-x_z),
 \qquad 1\leq z\leq n.
 \label{eq:compartments}
\end{equation}
Here each $f_z$ is an arbitrary measurable convex mixture
of the supplied endotactic fields with edge rates in
$[\kappa,K]$, and $m_{zw}$ are nonnegative measurable,
locally bounded coefficients. Each scalar transport
coefficient is the same for all species in its edge.
Then $\mathcal K^n$ is forward invariant, all solutions
starting there are global and positive, and in particular
every initial configuration with $x_z(0)\in Q$ stays in
this same range $\mathcal K$ at every compartment. The
one-compartment set does not depend on $n$, topology,
transport coefficient values, or switching paths.
\end{corollary}
\begin{proof}
At the fixed scale of Corollary~\ref{cor:safety} define the
finite affine minimum on product space
\[
 G(x_1,\ldots,x_n)=\min_{z,\lambda}L_{\lambda,h}(x_z)
                  =\min_z F_h(x_z).
\]
If $(z,\lambda)$ is active in this joint minimum, then
$\lambda$ is active at $x_z$, and for every compartment $w$,
\[
 L_{\lambda,h}(x_w)\geq F_h(x_w)\geq G
                         =L_{\lambda,h}(x_z).
\]
Hence
\[
 r_\lambda\cdot\sum_w m_{zw}(x_w-x_z)
   =\sum_w m_{zw}\bigl(L_{\lambda,h}(x_w)
                         -L_{\lambda,h}(x_z)\bigr)\geq0.
\]
The reaction contribution for this active label is also
nonnegative while $x_z\in B_h$, by the proof of
Corollary~\ref{cor:safety}. Lemma~\ref{lem:chain} on product
space therefore shows that $G$ is nondecreasing along every
solution segment in $B_h^n$.

Globally $G\leq M_h$, and $G=M_h$ exactly on
$\mathcal K_h^n$, since each $F_h\leq M_h$. Starting in that
product, monotonicity forces $G=M_h$ while in $B_h^n$.
The product plateau is compact and misses the boundary of
the box, so there is no first exit. Measurable locally
bounded transport and bounded reaction rates give local
Carath\'eodory uniqueness and continuation from this compact
positive set on every finite time interval. This proves
global positivity and invariance. The set's construction
preceded and did not involve the number or topology of
compartments.
\end{proof}

\begin{remark}
The transport assertion is invariance for a supplied compact
initial range. It does not assert a topology-independent
classwise absorber or convergence rate for arbitrary initial
configurations. Species-dependent transport matrices do not
satisfy the displayed affine-difference identity in general
and are not covered by this corollary.
\end{remark}

\section{Consequences, scope tests, and prior-art boundaries}
\label{sec:scope}

\begin{corollary}[Equilibria and periodically forced solutions]
\label{cor:periodic}
Under Theorem~\ref{thm:main}, fixed positive rates give a
positive equilibrium in each positive class. Measurable
$T$-periodic admissible rates give a positive $T$-periodic
Carath\'eodory solution in each class. The periodic conclusion
also holds under Theorem~\ref{thm:switching} when the mixture
weights and rates are $T$-periodic almost everywhere.
\end{corollary}
\begin{proof}
The relevant compact convex class set is invariant, and
uniform local state Lipschitz bounds imply continuity of
each solution map on it. For periodic data the time-$T$
map is a continuous selfmap, so Brouwer's theorem gives a
fixed point. Time periodicity and uniqueness imply that
the corresponding solution repeats with period $T$.

For autonomous fixed rates, apply Brouwer to the time-$1/n$
flow selfmap and choose a fixed point $z_n$ for each $n$.
On the compact invariant set the autonomous flow has the
uniform expansion $\Phi_t(z)=z+t f(z)+o(t)$ as $t\to0+$;
this follows by integrating $f(\Phi_u(z))$ and using uniform
continuity and boundedness of $f$. Thus $f(z_n)\to0$.
A convergent subsequence has a positive limit $z_*$ in the
compact set, and continuity gives $f(z_*)=0$.
\end{proof}

\paragraph{A strict extension example.}
In two species consider
\begin{equation}
 (1,3)\to(2,3),\qquad (3,3)\to(2,3),\qquad
 (0,0)\rightleftarrows(0,1).
 \label{eq:example}
\end{equation}
Each of its horizontal and vertical subnetworks is endotactic
in every direction. A finite union of endotactic diagrams is
endotactic: a globally maximal essential source is maximal
essential in any constituent diagram containing its edge,
so each such edge points inward. The horizontal arrows in
\eqref{eq:example} have no return paths, so it is not weakly
reversible. It is not strongly endotactic for $w=e_2$: the
overall maximal sources have height 3 and only horizontal,
$w$-neutral reactions, whereas all essential sources have
heights 0 or 1. Its stoichiometric span is $\R^2$.
Its constant-rate equations are
\[
 \dot x=y^3(k_1x-k_2x^3),\qquad \dot y=k_3-k_4y.
\]
For $r=-e_2$ the minimum essential tier is the source
$(0,1)$, whose negative vertical reaction gives positive
$r$-derivative. The overall minimum $r$-source tier instead
contains the neutral horizontal sources. This exposes the
importance of the essential-source definition in the proof.

\paragraph{A coupled three-species scope witness.}
The network
\begin{equation}
 2A\to2B\to2C\to2A,\qquad 4A\rightleftarrows6B,\qquad
 A+B\to2B
 \label{eq:3Dwitness}
\end{equation}
is endotactic. Its cycle and reversible pair form a weakly
reversible, hence endotactic, base diagram. To check the added
edge, write $w=(a,b,c)$. Its $w$-projection is $b-a$.
If this is positive, then $b>a$ and its source level $a+b$
is below $2b$. If $c\ne b$, the essential cycle edge
$2B\to2C$ has source level $2b$. If $c=b$, the essential
edge $2C\to2A$ instead has source level $2b$. Thus this
outward added edge is never on the maximal essential tier.
When its projection is negative it points inward, and when
it is zero it is not essential. The added edge therefore
preserves endotacticity. For completeness, a weakly
reversible diagram is endotactic because an outward edge
from a maximal essential source has a target at a strictly
higher level. A return path must contain a nonneutral edge;
the source of its first such edge is still at that higher
level, contradicting maximality among essential sources.

Three reaction vectors have determinant
\[
 \det\begin{pmatrix}
 -2&0&-4\\2&-2&6\\0&2&0
 \end{pmatrix}=8,
\]
so its stoichiometric span is $\R^3$. It is not weakly
reversible: no path returns from $2B$ to the source $A+B$.
It is not strongly endotactic for $w=-(1,1,1)$: every
overall maximal source has total degree 2 and is neutral
in this direction, while only the degree-4 and degree-6
reversible sources are essential. The full-span condition
ensures $w\notin S^\perp$. Hence this coupled higher-order
network lies outside the first-order, two-dimensional,
and strongly endotactic hypotheses directly discussed in
the cited prior results. Theorem~\ref{thm:main} gives a
common absorber on all of $\R_{>0}^3$ for its bounded
measurable rates; no claim about other possible realizations
or a prior classification of this individual example is made.

\paragraph{Exact scope falsifications.}
The supplied two-sided rate interval is necessary for the
stated conclusion. For the reversible endotactic network
$0\rightleftarrows A$, an inflow $1/(1+t)$ and outflow 1
give $x'=1/(1+t)-x$, and $x(t)\to0$ for every positive
initial value. One can see the limit by splitting its
variation-of-constants integral into early and recent
times. This violates persistence when the positive lower
rate bound is removed. Inflow $e^t$ and outflow 1 instead
give
\[
 x(t)=e^{-t}x(0)+(e^t-e^{-t})/2,
\]
which is unbounded when the upper rate bound is removed.
The one-way diagram $A\to2A$ is not endotactic and has
unbounded positive solutions. Remark~\ref{rem:span} gives
the separate counterexample to absorption under switching
between different spans. These are exact interface tests,
not numerical evidence for the general theorem.

\paragraph{State feedback and alternative kinetic exponents.}
The pointwise drift estimates apply to every absolutely
continuous trajectory whose velocity at almost every time is an
admissible rate-box or convex-mixture velocity. Thus they
also apply to all solutions of the corresponding differential
inclusion. An arbitrary discontinuous state-feedback law
need not admit an ordinary solution, so no such existence
or uniqueness claim is made. In systems where kinetic
exponents differ from stoichiometric sources, the relevant
diagram is the exponent-arrow diagram
$y_e\to y_e+\nu_e$. Endotacticity of an unrelated
stoichiometric graph is not a substitute for endotacticity
of this diagram. Delays, stochastic jump models, and
unrestricted external forcing are not consequences of the
deterministic theorem.

\paragraph{Prior art and what was supplied.}
Craciun--Nazarov--Pantea \cite{CNP2013} formulated the
extended variable-rate permanence conjecture and proved
the two-species endotactic case, also allowing power laws.
Gopalkrishnan--Miller--Shiu \cite{GMS2014} proved permanence
for strongly endotactic networks in arbitrary dimension;
their Conjecture 4.5 states the general endotactic target.
The toric-differential-inclusion embedding work
\cite{TDI2020} supplies geometric structure but an embedding
alone is not an absorbing-set proof. The recent
first-order result \cite{Xu2026} concerns a proper
subclass of arbitrary higher-order networks.
Release149 \cite{OAI149} supplies the unusually flexible
affine construction and the compact-plateau scale mechanism.
This manuscript attributes and fully reconstructs that
capital rather than claiming to have invented it. Its new
transfer is the weighted essential-source argument and its
common measurable-rate drift. A targeted live search and
the supplied release catalogue, inspected on 7 October
2026, did not identify another all-dimensional theorem
covering all endotactic measurable-rate diagrams. This
limited search is not a proof of priority.

\paragraph{Resources and limits of constructivity.}
The finite diagram, endotacticity assumption, class
representative, and global rate interval are supplied
inputs. No samples, retained rate history, equilibrium
oracle, or kinetic-balance identities are acquired by this
proof. The construction retains a finite affine family,
comparison vectors, its real slopes and offset functions,
and a chosen scale. Once the fixed-scale numerical coefficient
values are supplied or precomputed, evaluating its active labels
requires $O(|\Lambda|d)$ scalar operations. Evaluating or acquiring
the offset functions at that scale is separate preprocessing and
is not included in this online affine-evaluation count.
The proof establishes no useful bound on
$|\Lambda|$, numerical precision, the activity cutoff,
$h_P$, or the size of the waiting bound. Compactness
selection of a finite cover is not an implemented synthesis
algorithm or a polynomial-time certificate compiler.
An effective construction would have to charge all these
resources. Hash agreement of the pinned sources, agreement
between independent internal audits, and successful
typesetting do not establish mathematical correctness.

\begin{thebibliography}{9}
\bibitem{OAI149}
OpenAI, \emph{Uniform Permanence in Weakly Reversible
Mass-Action Systems}, 5 October 2026, release149.
Pinned repository commit
\texttt{adc7f1241b42e322a6451854ab7e4b4c146bf78a}.
Exact reconstructed source: \texttt{build/sections/affine.tex},
with plateau and absorption arguments in
\texttt{build/sections/trapping.tex} and
\texttt{build/sections/permanence.tex}.
\url{https://github.com/openai/math/tree/adc7f1241b42e322a6451854ab7e4b4c146bf78a/preprints/Uniform-Permanence-in-Weakly-Reversible-Mass-Action-Systems-October-5-2026}.

\bibitem{CNP2013}
G.~Craciun, F.~Nazarov, and C.~Pantea,
\emph{Persistence and permanence of mass-action and
power-law dynamical systems}, SIAM J. Appl. Math.
\textbf{73}(1) (2013), 305--329.
\url{https://doi.org/10.1137/100812355}.
Preprint: \url{https://arxiv.org/abs/1010.3050}.

\bibitem{GMS2014}
M.~Gopalkrishnan, E.~Miller, and A.~Shiu,
\emph{A geometric approach to the global attractor
conjecture}, SIAM J. Appl. Dyn. Syst. \textbf{13}(2)
(2014), 758--797.
\url{https://doi.org/10.1137/130928170}.
Definitions reconstructed from the exact v2 source:
\url{https://arxiv.org/html/1305.5303v2}.

\bibitem{TDI2020}
\emph{Endotactic Networks and Toric Differential Inclusions},
SIAM J. Appl. Dyn. Syst. (2020).
\url{https://doi.org/10.1137/19M1269312}.
Preprint: \url{https://arxiv.org/abs/1906.08384}.

\bibitem{Xu2026}
C.~Xu, \emph{First-order endotactic reaction networks},
J. Math. Biol. (2026).
\url{https://doi.org/10.1007/s00285-026-02378-2}.
Preprint: \url{https://arxiv.org/abs/2409.01598}.
\end{thebibliography}
\end{document}
